% Contoh hasil praktik Hari 2. Kompilasi dengan XeLaTeX.
\documentclass[aspectratio=169,11pt]{beamer}
\usepackage[indonesian]{babel}
\usepackage{booktabs,pgfplots}
\pgfplotsset{compat=1.18}
\usetheme{metropolis}
\setbeamertemplate{title page}{%
  {\usebeamerfont{title}\inserttitle\par}
  {\usebeamerfont{subtitle}\insertsubtitle\par}
  \usebeamertemplate{title separator}\medskip
  {\usebeamerfont{author}\insertauthor\par}
  {\usebeamerfont{date}\insertdate\par}}
\title{Jika panjang tali dua kali lipat, bagaimana periodenya?}
\author{Kelompok eksperimen ayunan}
\date{Pelatihan Beamer, 20 November 2026}

\begin{document}
\begin{frame}[plain]
  \titlepage
\end{frame}

\section{Model dan pengukuran}
\begin{frame}{Model memprediksi kenaikan sebesar faktor $\sqrt{2}$}
  Untuk simpangan kecil dan redaman lemah, periode memenuhi
  \[
    T \approx 2\pi\sqrt{L/g}.
  \]
  Panjang $L$ diukur dalam meter, periode $T$ dalam sekon,
  dan percepatan gravitasi $g$ dalam m/s$^2$.
  \uncover<2->{Pada $g$ yang sama, $T(2L)/T(L)\approx\sqrt{2}$.}
\end{frame}

\begin{frame}{Ukur panjang dari titik gantung ke pusat beban}
  \begin{columns}[c,onlytextwidth]
    \begin{column}{.43\textwidth}
      \centering
      \begin{tikzpicture}
        \draw[thick] (-1,0)--(1,0);
        \draw[dashed] (0,0)--(0,-2.7);
        \draw[thick] (0,0)--(-72:2.6) node[midway,right] {$L$};
        \fill (-72:2.6) circle (.14);
      \end{tikzpicture}
    \end{column}
    \begin{column}{.53\textwidth}
      Lepaskan beban dari simpangan kecil. Catat waktu $t$ untuk sepuluh
      getaran penuh. Satu periode adalah $T=t/10$.
    \end{column}
  \end{columns}
\end{frame}

\section{Hasil dan kesimpulan}
\begin{frame}{Periode naik dari 1,27 s menjadi 1,80 s}
  \centering
  \begin{tabular}{rrr}
    \toprule
    $L$ (m) & $t$ untuk 10 getaran (s) & $T$ (s)\\
    \midrule
    0,40 & 12,7 & 1,27\\
    0,60 & 15,6 & 1,56\\
    0,80 & 18,0 & 1,80\\
    \bottomrule
  \end{tabular}
  \par\bigskip Data ilustratif dari latihan Hari 1.
\end{frame}

\begin{frame}{Kuadrat periode bertambah sebanding dengan panjang tali}
  \centering
  \begin{tikzpicture}
    \begin{axis}[width=.73\linewidth,height=5cm,
      xlabel={$L$ (m)},ylabel={$T^2$ (s$^2$)},legend pos=north west]
      \addplot[only marks,mark=*] coordinates
        {(0.40,1.613) (0.60,2.434) (0.80,3.240)};
      \addplot[domain=.35:.85,no marks] {4.05*x};
      \legend{Data ilustratif,Garis $T^2=4{,}05L$}
    \end{axis}
  \end{tikzpicture}
  \par Gradien sekitar $4{,}05$ s$^2$/m memberi $g\approx9{,}75$ m/s$^2$.
\end{frame}

\begin{frame}{Menggandakan panjang memberi periode sekitar 1,42 kali}
  Rasio periode pada dua panjang ekstrem adalah
  \[
    \frac{1{,}80}{1{,}27}\approx1{,}42.
  \]
  Pada data ilustratif ini, hasilnya dekat dengan prediksi $\sqrt{2}$.
  Untuk menilai ketidakpastian pada eksperimen nyata, ulangi pengukuran
  waktu pada setiap panjang tali.
\end{frame}
\end{document}
